%% 自動生成（gen.py）
\bktocchap{4-10}{向きと大きさをもつ量：ベクトル}{2人で引くと、荷物はどっちへ動く？}{6}{2}{2}{340}{C}
\bktocchap{4-11}{ベクトルを成分で計算しよう}{川に流されながら渡る船は、どこに着く？}{6}{2}{3}{374}{C}
\bktocchap{4-12}{内積：どれだけ同じ向きか}{進む向きと直角に押しても、仕事は0}{6}{3}{4}{408}{C}
\bktocchap{4-13}{空間の座標と3次元のベクトル}{部屋の中の1点を、3つの数で表す}{6}{4}{4}{434}{C}
\bktocchap{4-14}{外積：回す力と右ねじ}{ドアのノブは、なぜちょうつがいから遠い？}{6}{4}{5}{468}{C}
\bktocchap{4-15}{いろいろな座標：極座標・円柱座標・球座標}{地球の上の場所は、緯度と経度で表す}{6}{4}{5}{500}{C}
